\section{The Bootloader Decides Your Memory Layout}
\label{sec:build-bootloader}

The single most dangerous line in an embedded project is the \codeinline{ORIGIN} of the flash region,
and it is usually copied from a tutorial for a board with no bootloader.

\subsection{What Is at the Base of Flash}

This board ships with a Gecko Bootloader occupying the first 24 KiB. On the EFR32MG24 flash begins at
\codeinline{0x08000000}, so:

\begin{table}[htbp]
    \centering
    \footnotesize
    \begin{tabular}{@{}p{0.26\textwidth}p{0.26\textwidth}p{0.42\textwidth}@{}}
        \toprule
        \textbf{Region} & \textbf{Addresses} & \textbf{Owner} \\
        \midrule
        Bootloader & \codeinline{0x08000000}--\codeinline{0x08005FFF} & factory; never written by this project \\
        Application & \codeinline{0x08006000}--end & every image in this book \\
        \bottomrule
    \end{tabular}
    \caption{24 KiB that a linker script must not touch.}
    \label{tab:flash-map}
\end{table}

\subsection{The Wrong memory.x, and What It Does}

The version that appeared in this project's first draft of the book — copied from a generic
\codeinline{cortex-m} example — read:

\begin{minted}{text}
MEMORY
{
  FLASH (rx) : ORIGIN = 0x00000000, LENGTH = 1536K
  RAM  (rwx) : ORIGIN = 0x20000000, LENGTH = 256K
}
\end{minted}

Every line is wrong for this board, and none of them produces an error. The linker happily lays the
image out from address zero, the programmer writes it exactly there, the write verifies — and the
board no longer boots, because the bootloader that was at zero has been replaced by an application
that expects to be at \codeinline{0x08006000}.

\subsection{The Version in Use}

\begin{minted}{text}
MEMORY
{
  /* The Gecko bootloader owns 0x08000000..0x08005FFF. An application that starts
     below 0x08006000 overwrites it. */
  FLASH (rx)  : ORIGIN = 0x08006000, LENGTH = 1536K - 24K
  /* The mailbox at 0x2003F000 is a convention shared by every example; the stack
     is placed below it so a deep call cannot scribble over the report. */
  RAM   (rwx) : ORIGIN = 0x20000000, LENGTH = 252K
}

_stack_start = ORIGIN(RAM) + LENGTH(RAM) - 4096;
\end{minted}

Both numbers are deliberate: 252K rather than 256K because the top 4 KiB are reserved, and
\codeinline{\_stack\_start} is below the mailbox so a stack overflow corrupts a report rather than the
data being reported.

\subsection{The Rule, Enforced in a Script}

A comment is not protection. The project has two mechanical guards:

\begin{minted}{bash}
# 1. where would this image land?
python3 scripts/elf-load-range.py target/thumbv8m.main-none-eabihf/release/examples/pac_led_blink
# load_min=0x08006000 load_max=0x08006523 size=1316 segments=3

# 2. refuse anything below the bootloader boundary, before writing
bash scripts/flash.sh "$ELF"          # dry run: prints the plan, refuses load_min < 0x08006000
bash scripts/flash.sh "$ELF" --confirm
\end{minted}

The first command is the one worth internalising: \textbf{read the load addresses before every flash},
not after a board stops answering. It takes one second and it is the difference between a mistake and
a brick.
